\worldchapter{Prägungen}{create}
\label{ch:marks}

\begin{multicols*}{2}
\section{Was Prägungen sind}\label{sec:was-pragungen-sind}\index{Prägung}

Prägungen\index{Prägung} sind die identitätsstiftenden Bausteine eines Charakters. Ein neuer Charakter hat drei Prägungen: eine \textbf{Abstammung}, einen \textbf{Weg} und eine \textbf{Bindung}.

Später kommt eventuell noch ein \textbf{Mal} hinzu, wenn Ereignisse bleibenden Eindruck beim Charakter hinterlassen haben.

Prägungen haben einen \textbf{Vorteil} und eine \textbf{Verwundbarkeit}.

\begin{itemize}[leftmargin=*]
\item Pro Prägung gilt immer genau der aktuelle Vorteils- und Verwundbarkeitssatz.
\item Wenn mehrere Vorteile denselben Zahlenbonus liefern, gilt nur der stärkste einzelne numerische Bonus.
\item Verwundbarkeiten sind keine optionalen Flufftexte.
Die Spielleiterin darf und soll sie ins Spiel ziehen.
\end{itemize}

\section{Bindung und Belastung}\label{sec:bindung-und-belastung}\index{Bindung}

Eine Bindung\index{Bindung} kann angerufen werden, wenn sie noch nicht belastet ist. Danach wird sie \textbf{belastet}.

Eine belastete Bindung gewährt keinen Bonus mehr, bis sie erneuert wird. Eine belastete Bindung kann auf zwei Arten \textbf{wiederhergestellt} werden:

\begin{itemize}[leftmargin=*]
\item durch eine bedeutende Szene der Versöhnung, Pflichterfüllung oder Fürsorge, die die Spielleiterin akzeptiert
\item durch eine erfolgreiche Auszeit-Aktion \textbf{Kontakt pflegen}
\end{itemize}

\section{Male}\label{sec:male2}\index{Mal}

Neue Charaktere starten ohne Mal, Male werden erst im Spiel erlangt.
Ein Charakter kann höchstens \textbf{3 Male} haben.

Für Male gilt:
\begin{itemize}[leftmargin=*]
\item die Verwundbarkeit ist immer aktiv
\item pro Probe darf nur ein Mal-Bonus wirken
\item Ruhe und normale Erholung entfernen kein Mal
\end{itemize}

\section{Wendepunkte und neue Male}\label{sec:wendepunkte-und-neue-male}
Wenn \textbf{Bürde}\index{Bürde} das Limit erreicht, steht der Charakter an einem \textbf{Wendepunkt}\index{Wendepunkt}.

Dann entscheidet die Spielerin:
\begin{itemize}[leftmargin=*]
\item sofort eine neues Mal nehmen
\item oder zusammenbrechen und die nächste Konfliktrunde aussetzen
\end{itemize}

Auch nach überlebtem \textbf{Sterben} muss ein neues Mal genommen werden.

\section{Erweiterte Prägung}\label{sec:erweiterte-pragung}
Wachstum kann eine zusätzliche Prägung freischalten.

Eine solche \textbf{Erweiterte Prägung}\index{Erweiterte Prägung} ist:
\begin{itemize}[leftmargin=*]
\item eng an eine Prägung gekoppelt
\item einmal pro Sitzung nutzbar
\item gewährt +10 oder senkt eine geringe Konsequenz
\item braucht immer eine passende Verwundbarkeit
\end{itemize}

\section{Praktische Auslegung}\label{sec:praktische-auslegung}
Ein Prägungs-Vorteil gilt immer dann, wenn seine Auslösebedingung in der Fiktion tatsächlich zutrifft.
Prägungen ersetzen keine Probe, sondern verändern Zielwert, Effekt oder zulässige Erzählungen.

\begin{ruleexample}
Ein Charakter mit dem Weg \textit{Spion} erhält +10 bei Infiltration und geheimer Kommunikation.
Derselbe Charakter bekommt daraus aber keinen Bonus, wenn sie offen vor Gericht um Gnade bittet.
Dort könnte eher eine passende Bindungs- oder Abstammungs-Regel greifen.
\end{ruleexample}

\section{Vollständiger Katalog}\label{sec:vollstandiger-katalog}
Alle Abstammungen, Berufungen, Bindungen und Male stehen gesammelt im Anhang \cref{app:marks}.
\end{multicols*}
